\documentclass[journal]{IEEEtran}
% IEEEtran forces \ttdefault to pcr (Courier); fall back to CM typewriter
% when Courier TFM files are not available in the TeX installation.
\renewcommand{\ttdefault}{cmtt}
\usepackage{amsmath,amssymb}
\usepackage{graphicx}
\usepackage{booktabs}
\usepackage{cite}
\usepackage{url}
\usepackage{microtype}
\usepackage{array}

\begin{document}

\title{RFSS: A 3GPP-Compliant Multi-Standard RF Signal Source Separation Dataset
with Realistic Channel and Hardware Impairments}

\author{Hao~Chen,~Rui~Jin,~and~Dayuan~Tan%
\thanks{H.~Chen, R.~Jin, and D.~Tan are with the Department of Electrical
Engineering. (\textit{Manuscript submitted 2026.})}%
}

\maketitle

% ============================================================
\begin{abstract}
% ============================================================
The coexistence of heterogeneous cellular standards (2G through 5G) in shared
spectrum has created complex electromagnetic environments that demand
sophisticated RF source separation techniques. However, no publicly available
dataset exists to support data-driven research on this problem. We present
RFSS (RF Signal Source Separation), a comprehensive open-source dataset of
100,000 multi-source RF signal samples generated with full 3GPP standards
compliance. The dataset covers four cellular standards --- GSM (3GPP TS~45.004),
UMTS (3GPP TS~25.211), LTE (3GPP TS~36.211), and 5G NR (3GPP TS~38.211) ---
with 2 to 4 simultaneous sources per sample (plus 4,000 single-source
reference samples in a companion file), at a common sample rate of
30.72~MHz. Each sample is processed through independent 3GPP TDL multipath
fading channels~\cite{3gpp38901} and realistic hardware impairments including
carrier frequency offset, I/Q imbalance, phase noise, DC offset, and PA
nonlinearity (Rapp model). Two complementary mixing modes are provided:
co-channel (all sources at baseband) and adjacent-channel (each source
frequency-shifted to its standard-specific carrier). The dataset totals 103~GB
in HDF5 format with a 70/15/15 train/validation/test split. We benchmark five
methods --- FastICA, beta-divergence NMF, Conv-TasNet~\cite{luo2019conv},
DPRNN~\cite{luo2020dual}, and a CNN-LSTM baseline --- using permutation-invariant
scale-invariant signal-to-noise ratio (PI-SI-SINR)~\cite{leroux2019sdr}.
Conv-TasNet achieves $-21.18$~dB PI-SI-SINR on 2-source mixtures versus
$-34.91$~dB for ICA, a 13.7~dB improvement. On co-channel mixtures, which
represent the primary separation challenge, Conv-TasNet reaches $-12.34$~dB
versus $-28.04$~dB for ICA and $-16.19$~dB for NMF. The dataset and
evaluation code are publicly available at
\url{https://huggingface.co/datasets/rfss/rfss-dataset}.
\end{abstract}

\begin{IEEEkeywords}
RF source separation, multi-standard dataset, 3GPP compliance, channel
modeling, hardware impairments, deep learning benchmark, permutation-invariant
training.
\end{IEEEkeywords}

% ============================================================
\section{Introduction}
% ============================================================

The proliferation of wireless communication technologies has produced
increasingly dense electromagnetic environments in which multiple cellular
generations coexist within shared or adjacent frequency
bands~\cite{cabric2004implementation,mitola1999cognitive}. Legacy GSM
infrastructure continues to carry machine-type traffic alongside LTE networks
that dominate mid-band spectrum, while 5G NR deployments overlap with both in
sub-6~GHz bands. This heterogeneous coexistence imposes persistent co-channel
and adjacent-channel interference that challenges conventional receivers
designed for single-standard operation.

RF source separation --- recovering the individual component signals from an
observed mixture --- is a natural response to this challenge. In contrast to
spectrum sensing~\cite{haykin2005cognitive,goldsmith2009breaking}, which
identifies the presence and type of signals, source separation seeks to
reconstruct each waveform so that downstream demodulation can proceed per
standard. Successful separation would enable flexible multi-standard receivers,
improve spectrum monitoring, and support interference coordination in cognitive
radio systems.

Machine learning has emerged as a powerful tool for physical-layer signal
processing~\cite{oshea2017introduction,west2017deep,oshea2018over}, with
deep architectures demonstrating compelling results on modulation classification
and channel estimation. Yet extending these successes to RF source separation
requires training data that accurately captures the diversity of multi-standard
mixtures under realistic propagation and hardware conditions. This data has not
previously existed.

Existing public RF datasets were designed for modulation recognition rather
than source separation. RadioML~\cite{oshea2016radio,oshea2016convolutional}
provides single-signal samples across synthetic modulation types;
it contains no multi-standard mixtures and no separation ground truth.
Other work focuses on single-channel characterization of individual
standards~\cite{rajendran2018deep,blossom2004gnu}, which is insufficient for
training or evaluating separation algorithms. Audio source separation
research has benefited enormously from curated datasets such as WSJ0-2mix and
MUSDB; the RF domain lacks an equivalent resource.

This paper addresses the gap with the following contributions.
\begin{enumerate}
  \item \textbf{Dataset scale and coverage.} 100,000 multi-source samples
    spanning all four major cellular generations, plus 4,000 single-source
    reference samples, totaling 103~GB. Every sample carries complete
    per-source ground-truth waveforms, enabling supervised training and
    rigorous evaluation.
  \item \textbf{3GPP-compliant signal generation.} GSM GMSK (TS~45.004),
    UMTS W-CDMA with Gold-code scrambling (TS~25.211), LTE OFDM with correct
    cyclic-prefix lengths (TS~36.211), and 5G NR with flexible numerology
    (TS~38.211) are all generated from first principles with parameters
    traceable to the cited 3GPP specifications.
  \item \textbf{Realistic channel and hardware impairments.} Independent
    3GPP TDL-A through TDL-E multipath fading channels and five categories
    of hardware impairment are applied per source before mixing.
  \item \textbf{Two mixing modes.} Co-channel and adjacent-channel scenarios
    capture distinct real-world interference regimes and enable separate
    performance analysis.
  \item \textbf{Benchmark suite.} Five methods spanning classical and
    deep-learning families are evaluated with permutation-invariant
    SI-SINR, establishing reproducible baselines for future work.
  \item \textbf{Open access.} Dataset, generation code, and evaluation
    scripts are publicly released.
\end{enumerate}

The remainder of the paper is organized as follows. Section~\ref{sec:related}
reviews related datasets and prior work. Section~\ref{sec:construction}
describes the dataset construction pipeline. Section~\ref{sec:characterization}
characterizes the resulting dataset. Section~\ref{sec:benchmark} presents
benchmark experiments and results. Section~\ref{sec:access} describes data
access and format. Section~\ref{sec:limitations} discusses limitations and
future work. Section~\ref{sec:conclusion} concludes.

% ============================================================
\section{Related Work}
\label{sec:related}
% ============================================================

\subsection{RF Datasets for Signal Intelligence}

The RadioML family of datasets~\cite{oshea2016radio,oshea2016convolutional,
oshea2018over} comprises single-signal recordings with varying modulation
type and SNR, and has become the standard benchmark for automated modulation
classification (AMC). While widely used, RadioML samples are not derived from
3GPP specifications: GSM, UMTS, LTE, and NR are absent as distinct
categories. More critically, RadioML contains only single-signal samples ---
there is no mixture of two or more simultaneous transmissions and therefore
no ground truth for separation.

DARPA's Spectrum Collaboration Challenge (SC2) produced operational spectrum
recordings across multiple waveforms, but these are oriented toward
resource-sharing protocol research rather than baseband signal recovery.
Work on distributed spectrum sensing~\cite{rajendran2018deep} similarly
targets classification rather than separation, and does not include the
full 2G--5G standard suite. GNU Radio toolkits~\cite{blossom2004gnu} allow flexible
waveform generation but provide no automated multi-standard mixing,
ground-truth labeling, or 3GPP compliance validation.

\subsection{Audio Source Separation Datasets}

The audio source separation community has established a strong tradition of
benchmark datasets that has directly informed the present work. WSJ0-2mix
introduced the two-speaker mixture paradigm and motivated permutation-invariant
training (PIT). WHAM! extended it with realistic noise. MUSDB provides
music-source separations with four stems. These datasets share a defining
property with RFSS: each mixture sample carries per-source ground-truth
references enabling supervised training with PIT loss, and evaluation with
SI-SNR~\cite{leroux2019sdr}. We adopt the same evaluation protocol for RF signals. The SI-SINR formula
used here is mathematically identical to the SI-SNR defined in
\cite{leroux2019sdr}; we adopt the SINR notation to emphasize that
interference, rather than additive noise, is the primary degradation in the
multi-standard RF coexistence setting.

The Conv-TasNet~\cite{luo2019conv} and DPRNN~\cite{luo2020dual} architectures
were originally developed on these audio benchmarks. We evaluate whether
they transfer to the RF domain without domain-specific modification, which
itself is a useful question for the community.

\subsection{RF Source Separation Prior Work}

Classical approaches to RF source separation rely on Independent Component
Analysis (ICA) or Non-negative Matrix Factorization (NMF).
ICA~\cite{hyvarinen2000ica} exploits statistical independence of source signals
and is the most commonly applied blind separation technique in spectrum
monitoring. NMF~\cite{fevotte2011nmf} decomposes the magnitude STFT into
spectral basis components and has been applied to overlapping narrowband
signals.

Neither approach benefits from learning across a large labeled corpus, which
limits their adaptability to the heterogeneous modulation types and channel
conditions encountered in 2G--5G coexistence. Deep learning methods for RF
source separation remain understudied, largely because no suitable training
dataset has been available. The RFSS dataset is designed to change this.

% ============================================================
\section{Dataset Construction}
\label{sec:construction}
% ============================================================

Figure~\ref{fig:pipeline} illustrates the complete dataset construction
pipeline. Each sample is produced by: (1) selecting 2--4 source standards
at random; (2) generating per-source baseband waveforms from 3GPP parameters;
(3) passing each waveform through an independent TDL channel and hardware
impairment chain; (4) mixing the impaired signals either co-channel or
adjacent-channel; and (5) storing the pre-impairment source waveforms as
separation targets alongside the mixture.

\begin{figure}[t]
  \centering
  \includegraphics[width=\columnwidth]{figures/fig_pipeline.pdf}
  \caption{RFSS dataset construction pipeline. Each sample passes through
  independent per-source signal generation, 3GPP TDL channel modeling,
  hardware impairment injection, and two-mode mixing before ground-truth
  targets and the mixture observation are written to HDF5.}
  \label{fig:pipeline}
\end{figure}

\subsection{Signal Generation}

All waveforms are generated at a common intermediate sample rate of 30.72~MHz,
the standard LTE/NR reference clock, before mixing. Each standard is
independently implemented from its 3GPP physical-layer specification.

\textbf{GSM (2G).} GMSK modulation per 3GPP TS~45.004~\cite{3gpp45004}:
\begin{equation}
\begin{split}
s_{\text{GSM}}(t) &= \exp\!\bigl(j\phi(t)\bigr), \\
\phi(t) &= \pi h \sum_k a_k \int_{-\infty}^{t} g(\tau - kT_s)\,\mathrm{d}\tau
\end{split}
\label{eq:gsm}
\end{equation}
where $h\!=\!0.5$ is the modulation index, $a_k \in \{-1,+1\}$ are
differentially encoded bits, $g(t)$ is the Gaussian pulse shaping filter with
bandwidth-time product $BT\!=\!0.3$, and $T_s\!=\!3.69\,\mu$s is the symbol
duration (270.833~ksps). The exponential form makes explicit that GMSK is a
constant-envelope modulation: signal energy enters only through the cumulative
phase $\phi(t)$, not through an amplitude envelope, yielding PAPR below 2~dB
and the 200~kHz channel bandwidth specified in TS~45.004.

\textbf{UMTS (3G).} W-CDMA per 3GPP TS~25.211~\cite{3gpp25211}:
\begin{equation}
s_{\text{UMTS}}(t) = \sum_{i=1}^{N_{\rm users}}
  \sqrt{P_i} \sum_k d_i[k]\, c_i[k]\, s_i[k]\; p(t - kT_c)
\label{eq:umts}
\end{equation}
where $P_i$ is per-user power, $d_i[k]$ are data symbols, $c_i[k]$ are OVSF
channelization codes, $s_i[k]$ is the Gold-code scrambling sequence per
TS~25.213, and $p(t)$ is the root-raised-cosine chip pulse with rolloff
$\alpha\!=\!0.22$ and chip duration $T_c\!=\!260.4$~ns (3.84~Mcps).
Channel bandwidth is 5~MHz.

\textbf{LTE (4G).} OFDM per 3GPP TS~36.211~\cite{3gpp36211}:
\begin{equation}
s_{\text{LTE}}(t) = \sum_{l=0}^{N_{\rm symb}-1}
  \sum_{k=0}^{N_{\rm SC}-1} X_l[k]\,
  e^{j2\pi k\Delta f(t - lT_s - T_{\rm CP})}
\label{eq:lte}
\end{equation}
where $X_l[k]$ are QAM symbols on subcarrier $k$ of OFDM symbol $l$,
$\Delta f\!=\!15$~kHz is subcarrier spacing, and $T_{\rm CP}$ is the cyclic
prefix duration as specified in TS~36.211 Table~6.12-1. The FFT size is
1024 for 10~MHz operation, yielding a nominal channel bandwidth of 10~MHz.

\textbf{5G NR (5G).} Flexible-numerology OFDM per 3GPP
TS~38.211~\cite{3gpp38211}:
\begin{equation}
s_{\rm NR}(t) = \sum_{l=0}^{N_{\rm symb}-1}
  \sum_{k=0}^{N_{\rm SC}-1} X_l[k]\,
  e^{j2\pi k\Delta f_{\rm SCS}(t - lT_s^\mu - T_{\rm CP}^\mu)}
\label{eq:nr}
\end{equation}
where $\Delta f_{\rm SCS}\!=\!15\!\times\!2^\mu$~kHz is the subcarrier
spacing for numerology $\mu \in \{0,1,2,3,4\}$, enabling bandwidths from
5~MHz to 400~MHz. The dataset uses $\mu\!=\!1$ (30~kHz SCS) with 50~MHz
nominal bandwidth and FFT size 2048, with cyclic-prefix durations per
TS~38.211 Table~5.3.1-2.

\subsection{Channel Modeling}

Each source signal passes through an independent single-input single-output
(SISO) channel drawn from the 3GPP TDL family defined in
TR~38.901~\cite{3gpp38901}: TDL-A and TDL-B are non-line-of-sight delay
profiles; TDL-C is the primary NLOS wideband profile; TDL-D and TDL-E are
Rician (line-of-sight) profiles with K-factors of 13.3~dB and 22~dB
respectively. Channel type is selected uniformly at random per source per
sample.

Time-varying Doppler is generated using Jakes' sum-of-sinusoids
model~\cite{jakes1994microwave}, with maximum Doppler frequency drawn
uniformly from 1 to 300~Hz. AWGN with SNR drawn uniformly from 0 to 30~dB
is added after channel filtering, providing a realistic range of received
signal quality.

The mixture model combining $N_{\rm src}$ impaired sources is:
\begin{equation}
y(t) = \sum_{i=1}^{N_{\rm src}} \sqrt{P_i}
  \sum_{l=0}^{L_i-1} h_{i,l}(t)\,
  s_i(t - \tau_i - \tau_{i,l})\,
  e^{j2\pi f_i t} + n(t)
\label{eq:mixture}
\end{equation}
where $P_i$ is the power scaling factor for source $i$, $h_{i,l}(t)$ is
the $l$-th tap of its TDL channel impulse response, $\tau_i$ is a
differential timing offset, $f_i$ is the carrier frequency shift
(zero for co-channel, standard-specific for adjacent-channel), and $n(t)$
is additive white Gaussian noise.

\subsection{Hardware Impairments}

Realistic receiver hardware impairments are modeled per 3GPP specifications
and applied independently to each source before
mixing~\cite{3gpp38104,3gpp38101,3gpp36101}:

\textit{Carrier Frequency Offset (CFO).} A multiplicative phase rotation
$e^{j2\pi f_{\rm CFO} t}$ with $f_{\rm CFO}$ drawn from $\pm0.05$--$5$~ppm
of the carrier, per TS~38.104 Section~6.5.1.

\textit{I/Q Imbalance.} Amplitude mismatch 0.1--3~dB and phase error
$1^\circ$--$10^\circ$, following the image rejection requirements of
TS~36.101 Section~7.9.

\textit{Phase Noise.} A Wiener-process phase perturbation with spectral
density $-90$ to $-110$~dBc/Hz at 10~kHz offset, per TS~25.102
Section~6.7.2.

\textit{DC Offset.} A constant complex additive offset at $-40$ to
$-30$~dBc relative to signal power, modeling local-oscillator leakage.

\textit{PA Nonlinearity.} The Rapp model~\cite{rapp1991effects} with
input back-off drawn from 3 to 9~dB below the amplifier saturation
point, producing realistic AM/AM and AM/PM distortion.

\subsection{Mixing Scenarios}

Two mixing modes are generated to capture distinct real-world interference
regimes.

\textbf{Co-channel mixing.} All sources are placed at baseband (no frequency
shift: $f_i\!=\!0$ for all $i$). Sources share the entire spectral region,
so separation must rely on differences in modulation type, temporal structure,
and the statistical properties of the learned channel. This is the most
challenging scenario and serves as the primary benchmark.

\textbf{Adjacent-channel mixing.} Each source is frequency-shifted by a
standard-specific offset before summing, so that the nominal carrier
frequencies of different standards do not overlap. This reduces spectral
overlap but introduces the frequency-management aspect of multi-standard
reception.

The number of simultaneous sources $N_{\rm src}$ is drawn from $\{2, 3, 4\}$
with probabilities reflecting realistic deployment densities. Across the
full 100,000-sample corpus, approximately 49\% of samples are two-source,
34\% are three-source, and 17\% are four-source. Each sample has a fixed
signal duration of 122,880 IQ samples at 30.72~MHz ($\approx 4$~ms).

% ============================================================
\section{Dataset Characterization}
\label{sec:characterization}
% ============================================================

\subsection{Composition}

The RFSS dataset comprises 100,000 multi-source samples partitioned into
training (samples 0--69,999; 70\%), validation (70,000--84,999; 15\%), and
test (85,000--99,999; 15\%) splits. An additional 4,000 single-source samples
are provided in a separate file (\texttt{rfss\_single.h5}) for signal
characterization and per-standard classification experiments.

Figure~\ref{fig:dataset_stats} shows the distribution of source counts,
mixing modes, and standard combinations across the corpus. Source count
follows an approximately geometric distribution (heavier weight on smaller
counts), and both mixing modes appear with roughly balanced frequency. All
$\binom{4}{k}$ standard combinations for $k \in \{2,3,4\}$ are represented.

\begin{figure}[t]
  \centering
  \includegraphics[width=\columnwidth]{figures/fig_dataset_stats.pdf}
  \caption{RFSS dataset composition statistics. Left: source-count
  distribution across 100,000 samples. Center: mixing-mode frequency
  (co-channel vs.\ adjacent-channel). Right: standard-combination
  frequency for multi-source samples.}
  \label{fig:dataset_stats}
\end{figure}

\subsection{Signal Properties}

Figure~\ref{fig:spectrograms} presents representative STFT spectrograms for
each of the four standards, illustrating the distinct time-frequency
signatures that a separator must disentangle. GSM exhibits a narrow-band
(200~kHz), constant-envelope character with quasi-stationary spectral
occupancy. UMTS spreads energy across 5~MHz with visible pseudo-noise
spreading gain. LTE and 5G NR display OFDM rectangular spectra with
subcarrier spacings of 15~kHz and 30~kHz respectively, and transient
variation across slots due to PDSCH scheduling.

\begin{figure}[t]
  \centering
  \includegraphics[width=\columnwidth]{figures/fig_spectrograms.pdf}
  \caption{Short-time Fourier transform spectrograms of single-source
  samples for all four cellular standards. Each panel shows time on the
  horizontal axis and frequency on the vertical axis, with power in dB
  encoded as color. The distinct spectral structures motivate the
  separability of multi-standard mixtures.}
  \label{fig:spectrograms}
\end{figure}

Figure~\ref{fig:signal_quality} characterizes three signal quality
dimensions. The PAPR panel confirms GSM's constant-envelope property
(PAPR $\approx\!1$--$2$~dB) versus LTE and 5G NR ($\approx\!11$--$13$~dB),
consistent with their multi-carrier waveforms. The power spectral density
panel highlights the 200-kHz GSM band, 5-MHz UMTS spread, 10-MHz LTE
rectangle, and 50-MHz NR rectangle. The amplitude distribution panel shows
the near-constant envelope of GSM contrasted with the Rayleigh-like tails
of LTE and 5G NR, as expected from the central-limit behavior of OFDM
symbols.

\begin{figure}[t]
  \centering
  \includegraphics[width=\columnwidth]{figures/fig_signal_quality.pdf}
  \caption{Signal quality characterization for the four cellular standards.
  Left: empirical PAPR distributions. Center: power spectral density
  estimates. Right: amplitude (envelope) probability density functions.}
  \label{fig:signal_quality}
\end{figure}

\subsection{HDF5 Format and Metadata}

The dataset is stored in two HDF5 files using gzip level-6 compression,
totaling 103~GB. The multi-source file (\texttt{rfss\_dataset.h5}) contains
four datasets:

\begin{itemize}
  \item \texttt{mixed\_signals} --- shape $(100000,\, 122880)$, complex64:
    the observed mixture waveform for each sample.
  \item \texttt{source\_signals} --- shape $(100000,\, 4,\, 122880)$,
    complex64: per-source ground-truth waveforms, zero-padded to a
    maximum of four sources.
  \item \texttt{signal\_lengths} --- shape $(100000,)$, int32: the active
    length of each source, supporting variable-length sources within the
    fixed array.
  \item \texttt{metadata} --- shape $(100000,)$, variable-length JSON
    string: per-sample metadata including source count, standard
    identifiers, mixing mode, SNR values, channel type, and hardware
    impairment parameters.
\end{itemize}

The flat array structure (no nested groups) simplifies HDF5 indexing for
random-access data loading, which is important for large-scale training.
All complex samples are stored as pairs of float32 (real, imaginary).

% ============================================================
\section{Benchmark Experiments}
\label{sec:benchmark}
% ============================================================

\subsection{Evaluation Metric}

We evaluate all methods with Permutation-Invariant Scale-Invariant
Signal-to-Interference-plus-Noise Ratio (PI-SI-SINR), following Le Roux et
al.~\cite{leroux2019sdr}. For a reference source $\mathbf{s}$ and an
estimated signal $\hat{\mathbf{s}}$, SI-SINR is:
\begin{equation}
\text{SI-SINR}(\hat{\mathbf{s}},\mathbf{s})
  = 10\log_{10}\frac{\|\alpha\mathbf{s}\|^2}
    {\|\hat{\mathbf{s}} - \alpha\mathbf{s}\|^2},\quad
  \alpha = \frac{\hat{\mathbf{s}}^\top\mathbf{s}}{\|\mathbf{s}\|^2}
\label{eq:sisinr}
\end{equation}
where mean subtraction is applied to both signals before the projection
to remove DC bias. Permutation invariance is achieved by selecting the
permutation of model outputs that maximizes the mean SI-SINR over all
sources.

The reported values are \emph{absolute} output PI-SI-SINR, not
improvement over input. Because the task is fully blind single-channel
separation (no spatial diversity), the input mixture SI-SINR is
effectively $-\infty$ (sources are completely overlapping), so absolute
output quality is the appropriate metric. Typical well-performing speech
separation systems achieve 8--15~dB absolute SI-SNR on WSJ0-2mix;
lower absolute values are expected here due to the increased difficulty
of the RF domain (unknown modulation type, channel distortion, hardware
impairments).

\subsection{Classical Baselines}

\textbf{FastICA.} We construct a Hankel (time-delay) embedding matrix
from the scalar observed mixture, increasing the effective number of
observations through lagged copies. FastICA then recovers statistically
independent components, which are matched to the reference sources by
permutation to maximize SI-SINR. ICA assumes statistical independence of
sources and non-Gaussianity; these assumptions hold approximately for
cellular signals but become weaker when sources share similar spectral
structure.

\textbf{Beta-divergence NMF.} The magnitude STFT of the mixture is
factored into $N_{\rm src}$ spectral bases and corresponding activations
using NMF with Kullback-Leibler divergence. Each source is then
reconstructed via Wiener-ratio masking in the STFT domain. NMF can
exploit differences in spectral shape (e.g., GSM's narrow band versus
LTE's wide rectangular spectrum) but cannot recover phase from the
magnitude spectrogram. Both baselines are evaluated permutation-invariantly
on $N\!=\!150$ test samples per source count.

\subsection{Deep Learning Methods}

Three architectures are trained from scratch on the RFSS training split.

\textbf{Conv-TasNet}~\cite{luo2019conv} uses a learned 1-D convolutional
encoder, a temporal convolutional network (TCN) with dilated depthwise
convolutions to compute source masks, and a learned decoder. We use
$N\!=\!256$ encoder filters, kernel length $L\!=\!16$, bottleneck dimension
$B\!=\!128$, skip channels $H\!=\!256$, kernel size $P\!=\!3$, $X\!=\!8$
blocks per repeat, and $R\!=\!3$ repeats.

\textbf{DPRNN}~\cite{luo2020dual} chunks the encoder output into
overlapping segments and applies alternating intra-chunk and inter-chunk
BiLSTM layers, enabling efficient global context modeling at low
computational cost. Configuration: $N\!=\!64$ encoder filters, $L\!=\!16$
kernel, feature dimension $B\!=\!64$, hidden dimension $H\!=\!64$,
chunk size $P\!=\!50$, 6 dual-path layers.

\textbf{CNN-LSTM} is our encoder-decoder baseline: a stack of dilated
causal convolutions extracts local features, a bidirectional LSTM layer
provides global sequence context, and a linear projection head estimates
each source waveform directly (regression, no masking). This architecture
makes the role of soft-mask nonlinearity explicit through comparison with
the two masking-based models above.

All models are trained with permutation-invariant training (PIT) loss
using negative SI-SINR summed over the best permutation of source estimates.
Optimization uses Adam with learning rate $10^{-3}$, gradient clip norm
1.0, batch size 8, and cosine annealing over 30 epochs with
$\eta_{\min}\!=\!10^{-5}$. Input crops of 7,680 samples (250~$\mu$s) are
randomly drawn during training. The best checkpoint per configuration
(lowest validation loss) is selected for evaluation on $N\!=\!300$ test
samples per source count. All experiments use single-channel (SISO) signals;
spatial diversity is not exploited.

\subsection{Results}

Table~\ref{tab:main} reports overall PI-SI-SINR on the test split.

\begin{table}[t]
\centering
\caption{Overall PI-SI-SINR (dB) on the test split
($N\!=\!150$ for baselines, $N\!=\!300$ for DL, seed\,42).
Higher values indicate better separation. Bold: best per row.}
\label{tab:main}
\resizebox{\columnwidth}{!}{%
\begin{tabular}{lrrrrr}
\toprule
Sources & ICA & NMF & CNN-LSTM & DPRNN & Conv-TasNet \\
\midrule
2-source & $-34.91$ & $-26.07$ & $-23.32$ & $-21.53$ & $\mathbf{-21.18}$ \\
3-source & $-36.98$ & $-29.69$ & $-23.65$ & $-21.31$ & $\mathbf{-21.08}$ \\
4-source & $-35.84$ & $-27.54$ & $-23.56$ & $-22.22$ & $\mathbf{-22.13}$ \\
\bottomrule
\end{tabular}}
\end{table}

Deep learning models consistently outperform classical baselines across all
source counts. Conv-TasNet and DPRNN are nearly tied (within 0.4~dB),
while CNN-LSTM trails by 1.4--2.4~dB. The improvement of Conv-TasNet over
ICA ranges from 13.7~dB (2-source) to 15.9~dB (3-source); over NMF, the
improvement is 4.9--8.6~dB across source counts. Performance degrades
modestly from 2-source to 4-source, consistent with the increased
combinatorial difficulty (24 permutations for PIT versus 2) and the smaller
per-configuration training set.\footnote{The Conv-TasNet 3-source result
($-21.08$~dB) uses a checkpoint from an earlier training run with
ReduceLROnPlateau scheduling, retained because it outperformed the
CosineAnnealingLR run on validation loss. All other reported checkpoints
used CosineAnnealingLR.}\footnote{The best DPRNN 4-source checkpoint
is from epoch~4 of 30. Validation loss stopped decreasing after this epoch,
likely due to the smaller 4-source training partition ($\approx\!10{,}500$
samples). The reported result reflects the best validation-selected
checkpoint, consistent with all other configurations.}

Figure~\ref{fig:benchmark} summarizes these results and additionally shows
the per-method breakdown by mixing mode.

\begin{figure}[t]
  \centering
  \includegraphics[width=\columnwidth]{figures/fig_benchmark.pdf}
  \caption{Benchmark results. Left: overall PI-SI-SINR (dB) for all methods
  across 2-, 3-, and 4-source configurations. Right: co-channel versus
  adjacent-channel PI-SI-SINR for all methods (2-source shown).
  Dashed boxes highlight the adjacent-channel evaluation floor
  (Section~\ref{sec:adjfloor}).}
  \label{fig:benchmark}
\end{figure}

Table~\ref{tab:co} reports co-channel PI-SI-SINR, which is the primary
measure of source separation capability.

\begin{table}[t]
\centering
\caption{Co-channel PI-SI-SINR (dB). $N_{\rm co}$ gives the number of
co-channel test samples per configuration. Bold: best per row.}
\label{tab:co}
\resizebox{\columnwidth}{!}{%
\begin{tabular}{lrrrrrr}
\toprule
Sources & $N_{\rm co}$ & ICA & NMF & CNN-LSTM & DPRNN & Conv-TasNet \\
\midrule
2-source & 110 & $-28.04$ & $-16.19$ & $-17.04$ & $-12.51$ & $\mathbf{-12.34}$ \\
3-source & 127 & $-28.20$ & $-15.08$ & $-15.99$ & $\mathbf{-10.38}$ & $-10.71$ \\
4-source & 133 & $-27.61$ & $-14.63$ & $-16.67$ & $-12.79$ & $\mathbf{-12.43}$ \\
\bottomrule
\end{tabular}}
\end{table}

On co-channel mixtures Conv-TasNet and DPRNN reach $-10$ to $-12$~dB versus
$-28$~dB for ICA, a 15--18~dB improvement. NMF performs at $-14$ to
$-16$~dB on co-channel, leveraging spectral shape differences between
standards. CNN-LSTM achieves $-15$ to $-17$~dB on co-channel, notably
closer to NMF than to Conv-TasNet/DPRNN, which underscores the importance
of the soft-mask architecture for this task.

\subsubsection{Adjacent-channel Evaluation Floor}
\label{sec:adjfloor}

The overall metric (Table~\ref{tab:main}) is substantially lower than the
co-channel metric (Table~\ref{tab:co}) for all methods. The reason is an
evaluation artifact, not a signal processing failure. Reference waveforms
in the dataset are stored in their pre-frequency-shift (baseband) form
for all sources. For adjacent-channel samples, the mixture contains
frequency-shifted versions of the sources; a separator that correctly
identifies and extracts a source still outputs a frequency-shifted estimate,
which the SI-SINR metric penalizes heavily when compared against the
baseband reference. This introduces a systematic floor of approximately $-28$~dB for
\emph{deep learning methods} on adjacent-channel samples. Classical baselines
(ICA: $-38$ to $-42$~dB on adjacent-channel) remain well below this floor,
reflecting both the frequency-shift penalty and genuine separation failure.

Approximately 37\% of test samples are co-channel and 63\% are
adjacent-channel, so the overall metric is dominated by this floor.
The co-channel breakdown (Table~\ref{tab:co}) is therefore the recommended
metric for comparing method capabilities. A future dataset release will
provide matched post-shift references to enable unambiguous evaluation
across both mixing modes.

% ============================================================
\section{Data Access and Format}
\label{sec:access}
% ============================================================

\subsection{Public Release}

The RFSS dataset is publicly available at:
\begin{center}
\url{https://huggingface.co/datasets/rfss/rfss-dataset}
\end{center}
The release includes the two HDF5 files (\texttt{rfss\_dataset.h5}, 103~GB;
\texttt{rfss\_single.h5}, 1.3~GB), along with the complete dataset generation
codebase, pre-trained model checkpoints, and evaluation scripts.

\subsection{HDF5 Access Example}

The following Python pseudocode illustrates loading a batch of samples:

\begin{verbatim}
import h5py, json
with h5py.File('rfss_dataset.h5', 'r') as f:
    x = f['mixed_signals'][idx]
    s = f['source_signals'][idx]
    n = f['signal_lengths'][idx]
    m = json.loads(f['metadata'][idx])
\end{verbatim}

The train/validation/test split is applied in code by index range:
training uses indices 0--69,999, validation 70,000--84,999, and test
85,000--99,999. No separate split files are needed.

\subsection{Evaluation Code}

A \texttt{SeparationDataset} PyTorch class is provided alongside
evaluation scripts for classical baselines and deep learning models.
These scripts reproduce all results in Section~\ref{sec:benchmark}
from the released checkpoints. A requirements file lists all Python
dependencies including PyTorch, h5py, scikit-learn, and scipy.

% ============================================================
\section{Limitations and Future Work}
\label{sec:limitations}
% ============================================================

\textbf{Adjacent-channel reference convention.} As described in
Section~\ref{sec:adjfloor}, storing pre-shift baseband references for
adjacent-channel samples creates an evaluation floor that obscures method
comparisons on those samples. The fix --- storing post-shift references as
well, or providing a flag for the applied frequency offset --- is
straightforward and will be included in the next dataset version.

\textbf{Single-antenna (SISO) only.} The current dataset models each
source as a scalar complex waveform with no spatial structure. Real
multi-antenna receivers can exploit spatial degrees of freedom for
separation; a MIMO extension (e.g., $2\times2$ or $4\times4$ arrays with
spatial correlation per 3GPP TDL) would unlock a richer class of
algorithms and closer correspondence to deployed base-station hardware.

\textbf{Downlink waveforms only.} The dataset covers downlink signals for
2G through 5G. Uplink waveforms differ in modulation order, power control
behavior, and resource allocation; their inclusion would broaden the
dataset's relevance to monitoring and coexistence analysis. Similarly,
NB-IoT, LTE-M, sidelink (D2D), and ISAC waveforms are absent.

\textbf{Future directions.} Planned extensions include: (1) incorporating
6G candidate waveforms (OTFS, filtered OFDM) as they approach
standardization; (2) using real-captured channel impulse responses (e.g.,
from COST or QuaDRiGa channel measurement campaigns) to replace the
synthetic TDL models; (3) scaling to higher source counts and wider
bandwidth mixtures; and (4) integrating automatic modulation classification
labels to support joint separation and classification tasks.

% ============================================================
\section{Conclusion}
\label{sec:conclusion}
% ============================================================

We presented RFSS, the first publicly available dataset specifically designed
for RF source separation research in heterogeneous cellular environments.
The dataset provides 100,000 labeled multi-source samples spanning GSM,
UMTS, LTE, and 5G NR with fully 3GPP-compliant signal generation, per-source
3GPP TDL fading channels, five categories of hardware impairment, and both
co-channel and adjacent-channel mixing modes, stored as 103~GB of HDF5 data.

Benchmarking with five methods demonstrates a clear performance hierarchy.
Conv-TasNet and DPRNN achieve $-10$ to $-12$~dB PI-SI-SINR on co-channel
mixtures, outperforming FastICA by 15--18~dB and beta-divergence NMF by
2--4~dB. CNN-LSTM, using direct waveform regression rather than masking,
trails the two masking-based architectures by 4--5~dB on co-channel
separation. All deep learning models substantially exceed classical baselines
across all configurations and mixing modes.

The RFSS dataset establishes reproducible baselines and provides the
training corpus needed to advance data-driven RF source separation research.
We release the dataset, generation code, trained models, and evaluation
scripts under open-source licenses to support the broader research community.

\bibliographystyle{ieeetr}
\bibliography{revised_paper}

\end{document}
